% !TEX program = xelatex
% DARTS-128 恢复路线报告 — 编译: xelatex 两次
\documentclass[UTF8,11pt,a4paper]{ctexart}

\usepackage{amsmath,amssymb,bm}
\usepackage{booktabs,array,tabularx,longtable}
\usepackage[margin=2.3cm]{geometry}
\usepackage{enumitem}
\usepackage{xcolor}
\setlist{itemsep=2pt,topsep=4pt}
\usepackage[colorlinks=true,linkcolor=blue!55!black,urlcolor=blue!55!black,citecolor=blue!55!black]{hyperref}
\usepackage{fancyhdr}

\setmainfont{Cambria}
\setmonofont{Consolas}
\newcommand{\code}[1]{\texttt{\small #1}}

\pagestyle{fancy}\fancyhf{}
\fancyhead[L]{\small DARTS-128（1--2b 模式）密钥恢复}
\fancyhead[R]{\small\thepage}
\renewcommand{\headrulewidth}{0.3pt}

\title{\vspace{-1.2cm}\textbf{DARTS-128（1--2b 模式）密钥恢复：\\公开可重放的双分支判定与 hint 通道}}
\author{GLM-5.3 Max \and JY Liu}
\date{2026 年 10 月 19 日}

\begin{document}
\maketitle
\thispagestyle{fancy}

\begin{abstract}\noindent
DARTS-128 是 NICCS 首轮征集的格签名候选。CryptHash 论坛其后发布了一次模式修订，把签名
响应中的共享符号因子从 $(1-b)$ 改为 $(1-2b)$。本文报告对该修订版参考实现的密钥恢复：
修订迫使验证侧采用 $\bar A z_1 \pm \hat q c\beta$ 双分支（验证方不知 $b$），而分支判定
可以被任何观察者公开重放——真分支恒匹配，假分支在 $\pm c\alpha$ 不跨量化边界时同样
匹配（概率 $(1-2^d/q)^{\tau} \approx 0.79$）。于是 200 万条公开签名中 21.0\% 携带
\textbf{零误差}的 $b$ 标签。以这些标签驱动两条独立的最小二乘通道（$\zeta$ 域标记最小
二乘与 hint 通道），再以公钥恒等式收尾，我们精确恢复 $(s_0, s_1, e)$ 全部 1536 字节，
与真值逐字节一致；未修改参考验证器下伪造 20/20 接受，三组阴性对照全部拒绝。完整技术
报告（14 页）见同目录 \code{DARTS\_1-2b\_攻击报告.pdf}。

\medskip\noindent\textbf{Abstract.}
DARTS-128 is a lattice signature candidate of the NICCS first round, later
revised by the CryptHash forum from a $(1-b)$ to a $(1-2b)$ shared sign factor.
The revision forces the verifier into a two-branch test $\bar A z_1 \pm \hat q
c\beta$, and this branch decision is publicly replayable: the true branch always
matches, while the false branch also matches whenever $\pm c\alpha$ does not
cross a quantization boundary (probability $(1-2^d/q)^{\tau}\approx 0.79$).
Consequently 21.0\% of 2 million public signatures carry a zero-error $b$ label.
Two independent least-squares channels driven by these labels---a $\zeta$-domain
tagged solve exploiting the negacyclic diagonalization, and the hint
channel---together with a public-key identity, recover $(s_0,s_1,e)$ exactly
(all 1536 bytes, byte-identical to the truth). Forgeries pass the unmodified
reference verifier 20/20; all three negative controls are rejected.
\end{abstract}

\section{对象与授权}

审计对象为 DARTS 提交包参考实现按 $(1-2b)$ 修订后的版本；修订涉及 \code{sign.c}
（三处）、\code{poly.c}、\code{poly.h} 与 \code{fft.c} 的一处整型修复，原文件以
\code{*.orig\_1b} 备份，逐文件 sha256 与批次清单在 \code{artifacts/}。研究在提交方授权
的自查场景下进行；攻击输入仅为公钥与公开签名，秘密密钥仅用于事后评分。评分与攻击严格
分离：\code{attack.py} 的真值参数为可选评分输入，攻击在完全不接触真值的模式下运行，
本报告全部攻击结果均在该模式下复现；\code{gen\_dump} 输出流的首字节为评分侧车（记录
真实 $b$），攻击代码不读取。

\section{泄露的来源：一个公开可重放的判定}

$(1-2b)$ 修订的本意是修补共享符号因子的分布缺陷，但它把验证方程变成了验证方对每个
候选符号的重放哈希：$\bar A z_1 + \hat q c\beta$ 或 $\bar A z_1 - \hat q c\beta$。
验证方不知道 $b$，只能两个都试；而``哪个分支匹配''这件事本身是公开信息。三类匹配
结果（仅 $+$、仅 $-$、双匹配、无匹配）中：

\begin{itemize}
\item \textbf{仅单侧匹配}的签名给出 $b$ 的确定标签。200 万条中 420{,}453 条（21.0\%），
  准确率 100.0000\%——不是统计意义上的高准确率，是\textbf{零误差}：单侧匹配在结构上
  排除另一个 $b$ 值。
\item 双匹配样本经生成测度的镜像对合 $(y,b,\mathrm{coins}) \mapsto (-y,1-b,\mathrm{coins})$
  严格证明 $P(b{=}0 \mid \text{双匹配}) = 1/2$，一阶统计无信息（实测
  49.9\%/50.1\%，$1.58\times10^6$ 条）——完整证明见 14 页报告~\S6。
\end{itemize}

\section{恢复路线}

\subsection{R1：$z_1$ 通道的 $\zeta$ 域标记最小二乘 $\to s_0, s_1$}

$z_1 = y_1 + (-1)^b (s_c \star c)$，$y_1$ 为超球新鲜噪声。标签已知时这是一个带符号的
线性回归，但 $R_q = \mathbb{Z}_q[x]/(x^n+1)$ 的乘法矩阵是\textbf{斜循环}的——直接套
循环自相关会使信号完全相消（我们实测过：真臂与硬币臂 rms 相同）。修复是精确的：负
循环在 $\zeta_k = e^{\,i\pi(2k+1)/n}$ 基下对角化，每条签名做 $2n$ 点 FFT、逐频点
除法，复杂度线性于签名数。\code{attack.py selftest} 内置与稠密最小二乘的自检（一致
到 $10^{-14}$）。

球拒绝采样带来一处必须处理的系统效应：$\|y\|$ 紧贴球半径（$d=1536$ 时标准差
$B/d \approx 3.9$），双球拒绝在 $\langle y, sc\rangle$ 方向构成近半空间截断，实测
$\mathbb{E}[y \mid \mathrm{acc}, b] = -0.63\cdot(-1)^b\cdot sc$——标记最小二乘估计
到的是 $\alpha = 0.37$ 倍缩放的信号。我们给出盲修正：矩估计
$\mathbb{E}[\hat\beta^2] = \alpha^2\cdot 0.3 + \sigma^2$（0.3 为三元秘密二阶矩，
$\sigma$ 取硬币臂），实测 $\hat\alpha = 0.367$，对真值 0.364。修正后加数据即可精确
四舍五入。200 万条下 $s_0$、$s_1$ 各 512/512 精确（rms 0.12）。

\subsection{R2 与 R2b：hint 通道与公钥恒等式 $\to e$}

hint 通道给出独立的第二条路线：
\[
h\cdot W \;\approx\; y_2 + (-1)^b (e \star c) - (1-2b)\,c,\qquad W = q/2^d,
\]
居中观测 $o = (-1)^b\cdot h\cdot W + c = (e \star c) + \text{对称噪声}$，同套标记最小
二乘直接估计 $e$。$h$ 的 mod-$2^d$ 环绕（\code{poly\_compress} 的掩码在 $q$ 边界附近
产生 $\pm 1020$ 编码符号）按模重居中即可还原真实小差值——编码器的边缘区间本为此设计，
对攻击者同样透明。

hint 通道先到 511/512；最后一处由公钥恒等式收尾：
\[
e \;=\; \hat q\,\beta - (A_0 s_0 + A_1 s_1).
\]
全解一致性 razor：$s$ 精确当且仅当 $e \in \{-1,0,1\}^{512}$。至此 $(s_0, s_1, e)$
全部 1536 字节与真值逐字节一致。

\subsection{伪造与阴性对照}

恢复密钥 $+$ 任取 $K$，未修改参考签名器/验证器下 20/20 接受。阴性对照三组：硬币符号
臂（退化噪声，无恢复）；篡改 $e$ / 篡改签名 / 错误 pk（全部拒绝）。伪造闭环与对照在
\code{forge.c} 中端到端。

\section{恢复程度}

\begin{table}[htbp]\centering\small
\begin{tabular}{@{}ll@{}}
\toprule
项 & 结果 \\
\midrule
干净标签 & 420{,}453 / 2{,}000{,}000（21.0\%），准确率 100.0000\% \\
$s_0$（$z_1$ 通道 LS） & 512/512 精确（rms 0.12） \\
$s_1$（$z_1$ 通道 LS） & 512/512 精确 \\
$e$（hint 通道 LS $+$ 公钥恒等式） & 511/512 $\to$ 恒等式修正 $\to$ 精确 \\
伪造 & 20/20 接受（未修改参考验证器） \\
阴性对照 & 三组全部拒绝 \\
\bottomrule
\end{tabular}
\end{table}

恢复的是\textbf{完整秘密钥}而非等价钥：DARTS 的密钥没有 BiT 式的验证等价自由度，
逐字节比对即验收。

\section{下一步工作}

\begin{enumerate}
\item \textbf{低数据量化}。当前 $2\times10^6$ 是充分条件；干净标签率 21\% 决定了标签
  噪声底，缩放修正已知（$\kappa \approx 0.63$ 的条件化收缩有盲修正），系统的
  $N$--精度曲线与自适应停止准则（\code{adaptive\_stop*.sh} 已有雏形）可以把数据量压到
  什么程度，是 \code{低数据攻击方案.md} 的主题，方案按数量级减小设计了三档。
\item \textbf{模糊签名的利用}。双匹配样本一阶无信息（\S2），但其条件分布仍依赖
  $(s,e)$ 与量化边界的相对位置——EM 式软标签在实测中发散（校准后 LLR 均值 0.7、后验
  翻转率 28\%，根因是拒绝采样条件化收缩同样压缩了模糊签名的双峰似然），这条路已被
  数据否定，记录在此防止重复投入。
\item \textbf{迁移}。任何``验证方必须枚举离散未知量''的修订都天然产生公开可重放判定；
  DARTS 案例的机理（分支匹配集而非单分支判定、斜循环域的正确对角化、球拒绝收缩的
  盲修正）对同类修订具有直接参考价值。
\end{enumerate}

\section{复现指南}

\begin{verbatim}
bash code/run_all.sh              # SMOKE=2000 全流程冒烟 / NS=500000 完整
python3 code/attack.py selftest   # zeta 域求解器 vs 稠密 LS 自检
\end{verbatim}

\code{artifacts/} 含攻击日志、批次清单与公开输入样本（pk、原始与重放 dump、恢复的
hint）；评分真值文件不入库（可由 \code{gen\_dump} 以记录的种子与批次参数再生，参数见
\code{artifacts/attack\_log.txt} 与 14 页报告~\S7）。修改后参考实现的构建流程与逐文件
sha256 见报告参考文献一节。

\section*{参考文献}
\begin{enumerate}[label={[\arabic*]}]
\item L.~Ducas, A.~Durmus, T.~Lepoint, V.~Lyubashevsky. \emph{Lattice Signatures and
  Bimodal Gaussians}. CRYPTO 2013（ePrint 2013/383）.
\item V.~Lyubashevsky. \emph{Fiat--Shamir with Aborts}. ASIACRYPT 2009; EUROCRYPT 2012.
\item Z.~Xiong, M.~Wang. \emph{Design-Level Cryptanalysis of the ICCS NGCC Round-1
  Public-Key Candidates}. ePrint 2026/2232.
\item DARTS 提交包：NICCS 首轮数字签名候选（DARTS-128 参数集，2026-09-20 归档版本）
  与 CryptHash 论坛 $(1-2b)$ 修订说明.
\end{enumerate}

\end{document}
